\chapter{Python Bot：同一个状态机如何在线运行}

\section{本章要解决的困惑}

fast backtest 读 \code{decision_frame}。Python Bot 在线运行时，不是直接读 parquet，而是从 Runner 的 public/private stream 接收消息，维护同样的 runtime-safe 状态。

这并不意味着 Bot 是另一套策略。它仍然是同一个状态机，只是输入形态从表格变成事件流。

\section{Bot 的职责}

Bot 只负责三件事：

\begin{enumerate}[leftmargin=2em]
  \item 从市场事件更新在线特征；
  \item 用 shadow policy 判断是否产生订单意图；
  \item 从 private fill/account 反馈更新闭合状态。
\end{enumerate}

它不能负责：

\begin{itemize}[leftmargin=2em]
  \item 控制 replay clock；
  \item 决定订单何时到达交易所；
  \item 计算 fill price；
  \item 读取未来 label；
  \item 读取 \code{date/} 研究输出。
\end{itemize}

\warningline{Bot 不读 date/。这不是风格偏好，而是防止 runtime 偷看研究结果的边界。}

\section{Python 文件阅读顺序}

\begin{lstlisting}
online_features.py
decision_frame.py
shadow_policy.py
execution_runtime.py
strategy.py
\end{lstlisting}

读法：

\begin{longtable}{p{0.28\textwidth}p{0.62\textwidth}}
\toprule
文件 & 先看什么 \\
\midrule
\code{online_features.py} & \code{OnlineFeatureState}、\code{RollingTrade}、\code{ClosedOutcomeState} \\
\code{decision_frame.py} & cache clock 如何从 trade/L2 build 一行状态 \\
\code{shadow_policy.py} & entry trigger、cell、gamma、pending entries \\
\code{execution_runtime.py} & admission、capacity、stopping helpers \\
\code{strategy.py} & stdin/stdout sparse Bot 怎样处理消息并输出 intent \\
\bottomrule
\end{longtable}

\section{online feature 和 decision frame 的关系}

两者都在描述“策略当前能看到什么”，但输入不同：

\begin{itemize}[leftmargin=2em]
  \item fast 线：已经有 \code{decision_frame_v1.parquet}，逐行读；
  \item Bot 线：收到 \code{market_quote}、\code{market_trade}、\code{market_l2_update}，边收边更新。
\end{itemize}

核心状态是一致的：

\begin{lstlisting}
last bid / ask / mid
spread_bps
rolling trades
trade_flow_imbalance
frames_since_mid_change
past_event_25_bps
closed outcome memory
\end{lstlisting}

\section{Bot 怎样得到 TFI}

\code{online_features.py} 里的 \code{add_trade} 会把每笔成交变成 signed quantity：

\begin{lstlisting}
buy  -> +qty
sell -> -qty
\end{lstlisting}

滚动窗口修剪后：

\[
TFI=\frac{signed\_qty}{abs\_qty+\epsilon}.
\]

如果窗口里没有成交，\code{abs_qty} 近似 0，TFI 返回 0。

\section{Bot 怎样得到 R5}

Bot 不能用未来收益。它只能从 private fill 或 fast shadow close 里得到已经闭合的 outcome。

\code{ClosedOutcomeState} 只保存最近 5 笔已闭合结果：

\begin{lstlisting}
outcomes_bps deque
positive_bps
negative_bps
r5
delta
energy
z
\end{lstlisting}

当前 runtime 只把 R5 用进 cell；Delta/Energy/Z 不进入控制。

\section{从 signal 到 order intent}

Bot 的 policy 产生的是一个策略意图，不是成交事实。意图里会有：

\begin{lstlisting}
side
target_exposure
feature_snapshot
observed_seq
observed_local_ts_us
\end{lstlisting}

Runner 收到以后，才决定：

\begin{itemize}[leftmargin=2em]
  \item 什么时候到达；
  \item 是否 accepted/rejected；
  \item 用 arrival quote 的 bid/ask 成交；
  \item portfolio 怎样变化。
\end{itemize}

\section{为什么 Bot 不该读 entries.parquet}

\code{entries.parquet} 是回测输出。它里面包含这笔 entry 是否被 capacity 接受、后来怎么退出等信息。Bot 如果读它，就相当于读了历史答案。

正确关系是：

\begin{lstlisting}
Bot state -> intent -> Runner/fill -> run artifacts
\end{lstlisting}

错误关系是：

\begin{lstlisting}
run artifacts -> Bot state -> intent
\end{lstlisting}

\checkpoint{Python Bot 是在线状态机。它和 fast backtest 的策略逻辑应尽量一致，但它不能读 fast 输出。}
